\documentclass[11pt,letterpaper]{article}
\usepackage[T1]{fontenc}
\usepackage[utf8]{inputenc}
\usepackage[scaled=0.95]{helvet}
\renewcommand{\familydefault}{\sfdefault}
\usepackage[top=1.7cm,bottom=1.5cm,left=2.2cm,right=2.2cm]{geometry}
\usepackage{microtype}
\usepackage[hidelinks]{hyperref}
\pagestyle{empty}
\setlength{\parindent}{0pt}
\setlength{\parskip}{7pt}

\begin{document}

{\large\bfseries Grover Aruquipa}\par\vspace{-6pt}
{\small Qu\'ebec, QC, Canada \ \textbullet\ \ grover.aruquipa.9@gmail.com \ \textbullet\ \ groveraruquipa.github.io}

\vspace{4pt}
Hiring Team, Promise Robotics, Calgary\\
Robotics Software Developer

Dear Hiring Team,

Robots that build houses in a factory have to do something most lab robots never do: run the same task thousands of times, on real material, without someone fixing them between cycles. I have spent my PhD making real robots work on the bench, and I would like to spend the next years writing the software that makes your production cells work every day.

I am finishing my PhD at Universit\'e Laval with Prof.\ Cl\'ement Gosselin and Prof.\ Patrice Lambert, and I defend in early 2027. Over three years I took three robots from the kinematic model to working hardware, and on all of them I wrote the software myself: the kinematics, the redundancy resolution, the collision and singularity checks, and the real-time control in C++ and MATLAB at 500\,Hz. I built a ROS~2 package with ros2\_control and MoveIt~2 planning for a 9-DoF robot, inside a Docker devcontainer, and I simulate my robots in MuJoCo and Isaac Sim before I run them on hardware. One of my papers under review is an analytical method to keep a robot's unlimited rotations free of collisions, which is exactly the kind of checking your motion planning needs in a tight cell.

Perception is not new to me either. Before my master's I worked as a computer vision engineer at Sencon, building industrial vision systems to detect scaffolding in industrial environments, and as a data science intern at Hitachi ABB Power Grids in Zurich, applying machine learning to high-voltage equipment. So I have already written code that had to work on a real site, with real lighting and real noise, not only on a dataset.

What I would bring to Promise is the habit of debugging where the problem actually is. On my robots I was the mechanical designer, the electrician and the programmer, so when a motion fails I can tell whether it is the planner, the model, the controller or a loose encoder. I write code that others can read, I am used to review from my supervisors and from journal referees, and I lead two graduate students day to day, so I know how to explain my work and take feedback quickly.

I will be eligible for an open Post-Graduation Work Permit as soon as I complete my PhD, so I do not need sponsorship, and I am ready to relocate to Alberta and work on site in Calgary and Edmonton from March 2027. Prof.\ Gosselin, Prof.\ Lambert and Univ.-Prof.\ Andreas M\"uller of JKU Linz can speak to my work; their details are in my CV.

Thank you for your time. I would be glad to walk you through my code and show the robots running.

Sincerely,\\
Grover Aruquipa

\end{document}
